% !PRINTED 227
\weberpage{227}
%<<CONT
the group nature has been proved and the icosahedral group has been formed.

From the first formula (25) there results further
\[
\psi = \chi\Theta^{\mp1}\chi\Theta^{\pm1}\chi\Theta^{\mp1},
\]
from which one concludes that $\psi$ can be derived from $\chi$ and $\Theta$, and hence that $\chi$ and $\Theta$ alone can already be regarded as \gesperrt{generating substitutions of the icosahedral group}.

If we wish to represent the substitution $\chi$ with the determinant 1, we observe the relation
\[
-\omega^2-1 = \omega-2 = (\varepsilon^2-\varepsilon^{-2})^2,
\]
and obtain, if we further substitute $\varepsilon+\varepsilon^{-1}$ for $\omega$:
\weq{26}{\chi = \begin{pmatrix} \dfrac{1}{\varepsilon-\varepsilon^{-1}}, & \dfrac{1}{\varepsilon^2-\varepsilon^{-2}} \\[2mm] \dfrac{1}{\varepsilon^2-\varepsilon^{-2}}, & \dfrac{-1}{\varepsilon-\varepsilon^{-1}} \end{pmatrix}.}

\S.\,59.

\gesperrt{The Divisors of the Icosahedral Group.}

The divisors of the icosahedral group must be sought among the lower polyhedral groups. Among these are found first the cyclic groups $C_n$, each of which consists of the powers of one of the icosahedral substitutions. The enumeration of the number of these groups follows at once from the arrangement of the substitutions according to their degrees, as we have given it in the preceding section; namely:
\[
\begin{aligned}
&15 \text{ groups } C_2, \\
&10 \quad \text{„} \quad C_3, \\
&\phantom{0}6 \quad \text{„} \quad C_5.
\end{aligned}
\]
There are further contained in the icosahedral group dihedral groups $D_2$, $D_3$, $D_5$, as whose representatives we set up the following:
\[
\begin{aligned}
D_2 &: 1, \psi, \chi, \psi\chi, \\
D_3 &: 1, \chi\Theta^2, \Theta^{-2}\chi, \psi\chi, \psi\Theta^2, \Theta^{-2}\psi\chi\Theta^2, \\
D_5 &: \Theta^r, \psi\Theta^r, \quad r=0,1,2,3,4,5.
\end{aligned}
\]
% [sic: Weber prints r=0,1,2,3,4,5; since \Theta^5=1, the last value duplicates r=0.]
